\section{Язык Lua}
\label{sec:lua}

\begin{infobox}[Цели главы]
\begin{itemize}
    \item Различать типы данных Lua и объявлять переменные --- локальные и глобальные.
    \item Записывать арифметику, сравнения и логические условия.
    \item Управлять ходом программы: ветвления \texttt{if}, циклы \texttt{for}, \texttt{while}, \texttt{repeat}.
    \item Писать свои функции и понимать, как автопилот вызывает нашу функцию \texttt{callback}.
    \item Узнавать типичные ошибки новичка по сообщению в терминале.
\end{itemize}
\end{infobox}

\subsection{Зачем Lua}

Lua --- лёгкий встраиваемый язык, созданный в начале 1990-х для расширения программ скриптами. Он компактен, не требует много памяти и быстро работает, поэтому живёт внутри игр, роботов и встраиваемых систем. Автопилот Пионера исполняет Lua прямо на борту: программа загружается в память контроллера и работает без компьютера.

В этом курсе Lua нужен, чтобы:
\begin{itemize}
    \item управлять светодиодами через модуль \texttt{Ledbar} (глава~\ref{sec:leds});
    \item выполнять действия по расписанию через \texttt{Timer}, не останавливая программу;
    \item отвечать на события автопилота --- взлёт завершён, точка достигнута, батарея садится --- в функции \texttt{callback(event)} (глава~\ref{sec:missions}).
\end{itemize}
Сейчас запоминать эти имена не нужно: в этой главе мы работаем только с языком.

\begin{rulebox}[Как работать с главой]
Каждый пример запускайте в симуляторе (язык \ui{Lua}, кнопка \ui{Запустить}) и сверяйте вывод в терминале с комментариями в коде: справа от \texttt{print} написано, что должно появиться.
\end{rulebox}

\subsection{Толковый словарь}
\begin{description}
  \item[\hypertarget{term:number}{Число}] вещественное или целое значение: \verb|3.14|, \verb|10|; математика --- в библиотеке \verb|math|.
  \item[\hypertarget{term:string}{Строка}] текст в кавычках: \verb|"Lua"|; склейка --- оператор \verb|..|, длина --- \verb|#s|.
  \item[\hypertarget{term:boolean}{Булево}] \verb|true| или \verb|false|; результат сравнения и материал для условий.
  \item[\hypertarget{term:nil}{nil}] «ничего»: переменная без значения, отсутствующее поле таблицы.
  \item[\hypertarget{term:table}{Таблица}] единственная составная структура Lua: массив \verb|t[1]| и словарь \verb|t.x| одновременно.
  \item[\hypertarget{term:function}{Функция}] именованный блок кода с параметрами и результатом.
  \item[\hypertarget{term:scope}{Область видимости}] \verb|local| делает переменную видимой только внутри текущего блока; без \verb|local| переменная глобальная --- видна во всём скрипте.
  \item[\hypertarget{term:truthiness}{Истинность}] ложными считаются только \verb|false| и \verb|nil|; ноль и пустая строка --- истина.
  \item[\hypertarget{term:ipairs}{ipairs / pairs}] обход таблицы: \verb|ipairs| --- по номерам 1..N, \verb|pairs| --- по всем ключам в произвольном порядке.
  \item[\hypertarget{term:unpack}{table.unpack}] раскладывает таблицу в список значений: \verb|table.unpack({1,0,0})| даёт \verb|1, 0, 0|. Часто пишут \verb|local unpack = table.unpack|.
\end{description}
